%! AP Statistics — Unit 2, Lesson 2.2 — Warm-Up
\documentclass[10pt]{article}
\usepackage{apstats-article}
\usepackage{apstats-boxes}

%=============================================================================
\begin{document}

\pageheader{Unit 2, Lesson 2.2}{Warm-Up}
\namedateperiod

\vspace{0.3in}

% ── SECTION 1: Spiral Review ──────────────────────────────────────────────────
{\large\bfseries\color{navy} Part 1 \quad Frequency Tables (Lesson 1.4 Review)}

\vspace{0.12in}

\noindent A survey asked 60 students which after-school activity they prefer.
The results are shown below.

\vspace{0.1in}

\renewcommand{\arraystretch}{1.6}
\begin{tabularx}{\linewidth}{@{} l C{2.8cm} C{2.8cm} @{}}
\rowcolor{sky}
\textbf{Activity}   & \textbf{Frequency} & \textbf{Relative Frequency} \\
\hline
Sports              & 24 & \blank{2.4cm} \\
Arts / Music        & 18 & \blank{2.4cm} \\
Academic clubs      & 12 & \blank{2.4cm} \\
None / Other        &  6 & \blank{2.4cm} \\
\hline
\textbf{Total}      & \blank{1.8cm} & \blank{2.4cm} \\
\end{tabularx}

\vspace{0.18in}

\begin{enumerate}

    \item Complete the \textbf{relative frequency} column. Show your work for the first row.\\[10pt]
    \writeline

    \vspace{0.1in}

    \item What percent of surveyed students preferred sports?\\[8pt]
    \writeline

    \vspace{0.1in}

    \item If you were to make a bar chart of this data, what would the height of the
    ``Arts / Music'' bar represent (in relative frequency form)?\\[8pt]
    \writeline

\end{enumerate}

\vspace{0.3in}

% ── SECTION 2: Looking Ahead ──────────────────────────────────────────────────
{\large\bfseries\color{navy} Part 2 \quad Two Variables at Once}

\vspace{0.12in}

\noindent Now suppose we also recorded whether each student is an underclassman
(9th/10th grade) or upperclassman (11th/12th grade).

\vspace{0.14in}

\begin{tcolorbox}[colback=greenbg, colframe=greenacc, boxrule=0.8pt, arc=2mm,
    left=4mm, right=4mm, top=3mm, bottom=3mm]
    \small
    \textbf{Think about it:} The one-way table above describes one variable at a time.
    If we want to know \emph{whether grade level is related to activity preference},
    we need to cross \textbf{two} categorical variables simultaneously.
    \begin{itemize}[leftmargin=1.3em, itemsep=6pt]
        \item A \textbf{two-way table} organizes two categorical variables at once
              --- you will build one in today's guided notes.
        \item What would the rows represent? What would the columns represent?
    \end{itemize}
\end{tcolorbox}

\vspace{0.16in}

\noindent\textbf{Sketch a blank two-way table} with grade level (Under / Upper)
as rows and activity (Sports / Arts / Clubs / None) as columns. Include a totals row and column.

\vspace{0.15in}

\renewcommand{\arraystretch}{2.2}
\begin{tabularx}{\linewidth}{@{} l Y Y Y Y Y @{}}
\rowcolor{sky}
\textbf{Grade} & \textbf{Sports} & \textbf{Arts} & \textbf{Clubs} & \textbf{None} & \textbf{Total} \\
\hline
Under          & & & & & \\
Upper          & & & & & \\
\hline
Total          & & & & & \\
\end{tabularx}

\end{document}
